\scanpage{7}
\mspage{5)}

\[
\operatorname{card}\alpha^{!-1}(\xi^{\unsure{\natural|\sharp|q|4}})=
\frac{\operatorname{card}D_{I_0}}{\operatorname{card}(G_\xi=\operatorname{Aut}(X,\Sigma,\Delta_{\unsure{*|i_0}}))}
=\frac{2n}{\deleted{\unsure{\mathrm{card}}}N(\xi^\natural)}
\]

Passant au composé

\diagram{\begin{tikzcd}[ampersand replacement=\&,row sep=large]
 \widetilde C_{\mathrm{st}}(S)\arrow[d,"\mathrm{deg.}\ 2"]\arrow[dd,bend left=50,"\widetilde\alpha"]\\
 C_{\mathrm{aff\,st}}^{!}(I_0)\arrow[d,"\alpha^!"]\\
 C_{\mathrm{aff}}(n)\simeq\widetilde C_{\mathrm{st}}(S)/D_S\simeq C_{\mathrm{aff\,st}}^{!}(I_0)/D_{I_0}
\end{tikzcd}}

on trouve

\[
\operatorname{card}\widetilde\alpha^{-1}(\xi^\natural)=\frac{4n}{N(\xi^\natural)}
\]

où

\[
N(\xi^\natural)=\operatorname{card}\operatorname{Aut}(X,\Sigma,\Delta)
\]

où

\[
\xi^\natural=\mathrm{cl}_{\mathrm{aff}}(X,\Sigma,\Delta)
\]

\uline{Ex} 1) $n=\operatorname{card}I_0=3$, alors $\operatorname{card}C_{\mathrm{aff}}(3)=1$, et $\operatorname{Aut}(X,\Sigma,\Delta)\simeq\mathfrak S_3$, de card $6$, donc $\widetilde C_{\mathrm{st}}(S)$ consiste en une seule orbite sous $D_{I_0}$, d'ordre $\frac{12}{6}=2$
% archive sheet number: 6
